\begin{table}[t]
\centering
\caption{Headline contrast family and test results. The family was assembled for transparent reporting rather than preregistered. H1--H4 use pooled benchmark-phase data; H5 uses the primary model.}
\label{tab:hypotheses}
\scriptsize
\setlength{\tabcolsep}{3pt}
\begin{tabular}{llllcccl}
\toprule
ID & Contrast & Test & Statistic & $N$ & $p$ & Effect & Outcome \\
\midrule
H1 & External-Benchmark Correlation & Pearson & $0.757$ & 1{,}800 & $<$0.001 & $0.757$ & Supported \\
H2 & Scramble Falsification & Pearson & $0.014$ & 1{,}200 & 0.638 & $0.014$ & Passes falsification \\
H3 & Directional Sensitivity & Pearson (1-sided) & $-0.791$ & 1{,}100 & $<$0.001 & $-0.791$ & Supported \\
H4 & Pooled Communication Contrast & Paired $t$ & $6.383$ & 1{,}479 & $<$0.001 & $0.166$ & Pooled effect; heterogeneous \\
H5 & Sender-Side Monitoring & Paired $t$ & $-19.109$ & 500 & $<$0.001 & $-0.855$ & Supported \\
\bottomrule
\end{tabular}
\begin{tablenotes}
\footnotesize\emph{Notes:} The Test column gives Pearson $r$ for H1--H3 and paired $t$-tests for H4--H5. H4 is a cell-pooled contrast whose equal-weighted model average is near zero, so its outcome records heterogeneity. H5 uses the primary model and matched prompt-isolation cells. The ex-post Bonferroni calculation is descriptive, not confirmatory. Survivors at $\alpha = 0.01$: H1, H3, H4, H5.
\end{tablenotes}
\end{table}
